% problems2.tex — Problem Set and Answers for Chapter 2

\begin{problemset}

\subsection*{Analytical Exercises}

\exercise[1] (Taken from Example 3.4.2 of Amemiya (1985))\quad
Let $z_n$ be defined by
\[
z_n = \begin{cases}
0 & \text{with probability } (n-1)/n, \\
n^2 & \text{with probability } 1/n.
\end{cases}
\]
Show that $\plim_{n\to\infty} z_n = 0$ but $\lim_{n\to\infty} \E(z_n) = \infty$.

\exercise[2] Prove Chebychev's weak LLN, by showing that
$\bar{z}_n \to_{\text{m.s.}} \mu$. \textbf{Hint:}
\[
\bar{z}_n - \mu = (\bar{z}_n - \E(\bar{z}_n)) + (\E(\bar{z}_n) - \mu).
\]
So
\[
(\bar{z}_n - \mu)^2 = (\bar{z}_n - \E(\bar{z}_n))^2
+ 2(\bar{z}_n - \E(\bar{z}_n))(\E(\bar{z}_n) - \mu)
+ (\E(\bar{z}_n) - \mu)^2.
\]
Show that $\E\!\big[(\bar{z}_n - \E(\bar{z}_n))(\E(\bar{z}_n) - \mu)\big] = 0$.

\exercise[3] (Consistency and Asymptotic Normality of OLS for Random Samples)\quad
Consider replacing Assumption~2.2 by

\medskip
\noindent\textbf{Assumption 2.2$'$:}\quad $\{y_i, \mathbf{x}_i\}$ \emph{is a random sample.}

\medskip
\noindent and Assumption~2.5 by

\medskip
\noindent\textbf{Assumption 2.5$'$:}\quad
$\mathbf{S} \equiv \E(\mathbf{g}_i\mathbf{g}_i')$ \emph{exists and is finite.}

\medskip
\noindent
Prove the following simplified version of Proposition~2.1:

\medskip
\noindent\textbf{Proposition~2.1 for Random Samples:}\quad
\emph{Under Assumptions~2.1, 2.3, 2.4, and 2.2$'$, the OLS estimator $\mathbf{b}$ is consistent.
Under the additional assumption of Assumption~2.5$'$, the OLS estimator $\mathbf{b}$ is
asymptotically normal with $\avar(\mathbf{b})$ given by (2.3.4).}

\begin{enumerate}[label=(\alph*)]
\item Show that this is a special case of Proposition~2.1 of the text.
\item Give a direct proof of this proposition.
\textbf{Hint:} Use Kolmogorov and Linderberg--Levy.
\end{enumerate}

\exercise[4] (Proof of Proposition~2.4)\quad
We wish to prove Proposition~2.4. To avoid inessential complications, we assume as in the text that $K = 1$ (only one regressor).
What remains to be shown is that the sample mean of the first term on the RHS
of (2.5.3) converges in probability to some finite number. Prove it.
\textbf{Hint:} Let $f \equiv x_i\varepsilon_i$ and $h \equiv x_i^2$. The \textbf{Cauchy--Schwartz inequality} states:
\[
\E(|f \cdot h|) \leq \sqrt{\E(f^2)\,\E(h^2)}.
\]
So
\[
\E(|x_i^3\varepsilon_i|) \leq \sqrt{\E(x_i^2\varepsilon_i^2)\,\E(x_i^4)}.
\]
By assumption, $\E(x_i^2\varepsilon_i^2)$ ($= \E(g_i^2)$) is finite, and by Assumption~2.6
$\E(x_i^4)$ is also finite.
So $\E(\varepsilon_i x_i^3)$ is finite and
$\frac{1}{n}\sum_i \varepsilon_i x_i^3 \pto$ some finite number by ergodic stationarity.

\exercise[5] (Direct proof that the change in $SSR$ divided by $\sigma^2$ is asymptotically
$\chi^2$)\quad
In Section~2.6, we proved that
\[
(SSR_R - SSR_U)/s^2 \dto \chi^2(\#\mathbf{r})
\]
as a corollary to Proposition~2.3. Give a direct proof, first by showing that
\[
SSR_R - SSR_U = (\sqrt{n}\,\bar{\mathbf{g}})'
\mathbf{S}_{\mathbf{xx}}^{-1}\,\mathbf{R}'
(\mathbf{R}\,\mathbf{S}_{\mathbf{xx}}^{-1}\,\mathbf{R}')^{-1}\,
\mathbf{R}\,\mathbf{S}_{\mathbf{xx}}^{-1}\,(\sqrt{n}\,\bar{\mathbf{g}}),
\]
\[
\bar{\mathbf{g}} \equiv \frac{1}{n}\sum_{i=1}^{n}\mathbf{x}_i \cdot \varepsilon_i.
\]

\exercise[6] (Optional, consistency of $\hat{\bm{\alpha}}$)\quad
In this question we prove the claim made in Section~2.8 that $\hat{\bm{\alpha}}$, the OLS estimate for the regression of $e_i^2$ on $\mathbf{z}_i$, is consistent.

Let $\tilde{\bm{\alpha}}$ be the OLS estimate from
%
\begin{equation}
\varepsilon_i^2 = \mathbf{z}_i'\bm{\alpha} + \eta_i,
\quad \eta_i \equiv \varepsilon_i^2 - \E(\varepsilon_i^2 \mid \mathbf{x}_i).
\tag{2.8.8}
\end{equation}

\begin{enumerate}[label=(\alph*)]
\item Make appropriate assumptions in addition to Assumptions~2.1 and 2.2 to
show that $\tilde{\bm{\alpha}}$ is consistent.
\textbf{Hint:} You need a rank condition for $\mathbf{z}_i$.

\item Derive:
\[
\hat{\bm{\alpha}} - \tilde{\bm{\alpha}}
= \Bigl(\frac{1}{n}\sum_{i=1}^{n}\mathbf{z}_i\mathbf{z}_i'\Bigr)^{-1}
\frac{1}{n}\sum_{i=1}^{n}\mathbf{z}_i \cdot v_i
\tag{$*$}
\]
with $v_i \equiv -2(\mathbf{b} - \bm{\beta})'\mathbf{x}_i\cdot\varepsilon_i
+ (\mathbf{b} - \bm{\beta})'\mathbf{x}_i\mathbf{x}_i'(\mathbf{b} - \bm{\beta})$.
\textbf{Hint:} The discrepancy between $\varepsilon_i^2$ and the squared OLS residual $e_i^2$ from the
original equation $y_i = \mathbf{x}_i'\bm{\beta} + \varepsilon_i$ is given by (2.3.10). Substituting it into
(2.8.8) gives: $e_i^2 = \mathbf{z}_i'\bm{\alpha} + (\eta_i + v_i)$.

\item To avoid inessential complications, suppose that $x_i$ is a scalar $x_i$. Then
$\frac{1}{n}\sum_i \mathbf{z}_i \cdot v_i$ in ($*$) becomes
\[
-2(b - \beta)\,\frac{1}{n}\sum_{i=1}^{n}x_i\varepsilon_i\mathbf{z}_i
+ (b - \beta)^2\,\frac{1}{n}\sum_{i=1}^{n}x_i^2\mathbf{z}_i.
\tag{$**$}
\]
Show that the plim of the first term is zero. \textbf{Hint:}
\[
\E(x_i\varepsilon_i\mathbf{z}_i) = \E[\mathbf{z}_i \cdot x_i \cdot \E(\varepsilon_i \mid x_i)].
\]
Show that the plim of the second term is zero if $\E(x_i^2\mathbf{z}_i)$ exists and is finite.

\item Thus we have proved that $\hat{\bm{\alpha}} - \tilde{\bm{\alpha}}$ vanishes. Prove the stronger result that
$\sqrt{n}(\hat{\bm{\alpha}} - \tilde{\bm{\alpha}})$ vanishes. \textbf{Hint:} Use Lemma~2.4(b).
\end{enumerate}

\exercise[7] (Optional, proof that WLS is asymptotically more efficient than OLS)\quad
In Section~2.8, the WLS estimator is denoted $\hat{\bm{\beta}}(\widetilde{\mathbf{V}})$.
We wish to prove that it is asymptotically more efficient than OLS, namely,
\[
\avar(\hat{\bm{\beta}}(\widetilde{\mathbf{V}})) \leq \avar(\mathbf{b}).
\]
But since $\avar(\hat{\bm{\beta}}(\widetilde{\mathbf{V}})) = \avar(\hat{\bm{\beta}}(\mathbf{V}))$,
it is sufficient to prove that
\[
\avar(\hat{\bm{\beta}}(\mathbf{V})) \leq \avar(\mathbf{b}).
\]
Prove this last matrix inequality. \textbf{Hint:} In Chapter~1, we proved the algebraic
result that
\[
(\mathbf{X}'\mathbf{X})^{-1}(\mathbf{X}'\mathbf{V}\mathbf{X})(\mathbf{X}'\mathbf{X})^{-1}
\geq (\mathbf{X}'\mathbf{V}^{-1}\mathbf{X})^{-1}
\tag{$*$}
\]
for any positive definite matrix $\mathbf{V}$. (This follows from the fact that GLS is more
efficient than OLS; the LHS is ($\sigma^2$ times) the variance of the OLS estimator for
the generalized regression model while the RHS is ($\sigma^2$ times) the variance of
the GLS estimator.) Dividing both sides of ($*$) by $n$ and taking probability limits
yields:
\[
(\plim \tfrac{1}{n}\mathbf{X}'\mathbf{X})^{-1}
(\plim \tfrac{1}{n}\mathbf{X}'\mathbf{V}\mathbf{X})
(\plim \tfrac{1}{n}\mathbf{X}'\mathbf{X})^{-1}
\geq \plim(\tfrac{1}{n}\mathbf{X}'\mathbf{V}^{-1}\mathbf{X})^{-1}.
\tag{$**$}
\]
The RHS is $\avar(\hat{\bm{\beta}}(\mathbf{V}))$ (see (2.8.7)). Show that the LHS equals
$\avar(\mathbf{b}) = \bm{\Sigma}_{\mathbf{xx}}^{-1}\,\mathbf{S}\,\bm{\Sigma}_{\mathbf{xx}}^{-1}$.

\exercise[8] Prove \eqref{eq:2.9.7}. \textbf{Hint:} If $(\mu, \bm{\gamma})$ is the least squares projection
coefficients, it satisfies
\[
\E(\mathbf{x}\mathbf{x}')
\begin{bmatrix} \mu \\ \bm{\gamma} \end{bmatrix}
= \E(\mathbf{x} \cdot y),
\tag{$*$}
\]
or
\[
\begin{bmatrix} \mu \\ \bm{\gamma} \end{bmatrix}
= [\E(\mathbf{x}\mathbf{x}')]^{-1}\,\E(\mathbf{x} \cdot y),
\tag{$**$}
\]
where
\[
\E(\mathbf{x}\mathbf{x}') = \begin{bmatrix}
1 & \E(\widetilde{\mathbf{x}})' \\
\E(\widetilde{\mathbf{x}}) & \E(\widetilde{\mathbf{x}}\widetilde{\mathbf{x}}')
\end{bmatrix}, \quad
\E(\mathbf{x} \cdot y) = \begin{bmatrix}
\E(y) \\ \E(\widetilde{\mathbf{x}} \cdot y)
\end{bmatrix}.
\]
The first equation of ($*$) is
$\mu + \E(\widetilde{\mathbf{x}})'\bm{\gamma} = \E(y)$. Use this to eliminate $\mu$ from the
rest of the equations of ($*$) and then solve for $\bm{\gamma}$. In the process, use
$\Var(\widetilde{\mathbf{x}}) = \E(\widetilde{\mathbf{x}}\widetilde{\mathbf{x}}') - \E(\widetilde{\mathbf{x}})\,\E(\widetilde{\mathbf{x}})'$,
$\Cov(\widetilde{\mathbf{x}}, y) = \E(\widetilde{\mathbf{x}} \cdot y) - \E(\widetilde{\mathbf{x}})\,\E(y)$.
Alternatively, use the formula for the partitioned inverse (see (A.10) of Appendix~A) to obtain
\[
\E(\mathbf{x}\mathbf{x}')^{-1} = \begin{bmatrix}
1 + \E(\widetilde{\mathbf{x}})'\,\Var(\widetilde{\mathbf{x}})^{-1}\,\E(\widetilde{\mathbf{x}})
& -\E(\widetilde{\mathbf{x}})'\,\Var(\widetilde{\mathbf{x}})^{-1} \\
-\Var(\widetilde{\mathbf{x}})^{-1}\,\E(\widetilde{\mathbf{x}})
& \Var(\widetilde{\mathbf{x}})^{-1}
\end{bmatrix}.
\]
Then substitute this into ($**$).

\exercise[9] (Proof of Proposition~2.9 for $p = 1$)\quad
Assume that $\{\varepsilon_t\}$ is an ergodic stationary
martingale difference sequence and that
$\E(\varepsilon_t^2 \mid \varepsilon_{t-1}, \varepsilon_{t-2}, \ldots, \varepsilon_1) = \sigma^2 < \infty$.
Let $g_t = \varepsilon_t \cdot \varepsilon_{t-1}$ and
\[
\hat{\gamma}_j = \frac{1}{n}\sum_{t=j+1}^{n}\varepsilon_t \cdot \varepsilon_{t-j}
\quad (j = 0, 1), \quad
\hat{\rho}_1 = \frac{\hat{\gamma}_1}{\hat{\gamma}_0}.
\]

\begin{enumerate}[label=(\alph*)]
\item Show that $\{g_t\}$ ($t = 2, 3, \ldots$) is a martingale difference sequence.
\item Show that $\E(g_t^2) = \sigma^4$.
\item Show that $\sqrt{n}\,\hat{\gamma}_1 \dto N(0, \sigma^4)$ as $n \to \infty$.
\item Show that $\sqrt{n}\,\hat{\rho}_1 \dto N(0, 1)$.
\textbf{Hint:} First show that
$\sqrt{n}\,\hat{\gamma}_1/\sigma^2 \dto N(0, 1)$. Then use Lemma~2.4(c).
\end{enumerate}

\exercise[10] (Asymptotics of sample mean of MA(2))\quad
Let $(\varepsilon_{-1}, \varepsilon_0, \varepsilon_1, \varepsilon_2, \ldots)$ be i.i.d.\
with mean zero and variance $\sigma_\varepsilon^2$. Consider a process
$(y_{-1}, y_0, y_1, y_2, \ldots)$ generated by
\[
y_{-1} = \varepsilon_{-1}, \quad
y_0 = \varepsilon_0 + \theta_1\varepsilon_{-1}, \quad
y_t = \varepsilon_t + \theta_1\varepsilon_{t-1} + \theta_2\varepsilon_{t-2}
\quad (t = 1, 2, \ldots\,).
\]
(This process is called a second-order moving average process (MA(2)); see
Chapter~6.)
\begin{enumerate}[label=(\alph*)]
\item Show that $(y_1, y_2, \ldots)$ is covariance stationary. Derive the expressions for
the autocovariances $\gamma_j$ ($j = 0, 1, 2, \ldots$).

\item Let
\begin{align*}
r_{tj} &\equiv \E(y_t \mid y_{t-j}, y_{t-j-1}, \ldots, y_0, y_{-1}) \\
&\quad - \E(y_t \mid y_{t-j-1}, y_{t-j-2}, \ldots, y_0, y_{-1}) \\
&\quad (t = j, j+1, \ldots;\; j = 0, 1, 2, \ldots\,).
\end{align*}
Show that
\[
r_{t0} = \varepsilon_t, \quad
r_{t1} = \theta_1\varepsilon_{t-1}, \quad
r_{t2} = \theta_2\varepsilon_{t-2}, \quad
r_{t3} = 0, \quad
r_{t4} = 0, \quad \ldots
\]
\textbf{Hint:} There is a one-to-one mapping between
$(\varepsilon_{-1}, \varepsilon_0, \varepsilon_1, \ldots, \varepsilon_t)$ and
$(y_{-1}, y_0, y_1, \ldots, y_t)$ so
$\E(y_t \mid y_{t-j}, \ldots, y_{-1}) = \E(y_t \mid \varepsilon_{t-j}, \ldots, \varepsilon_{-1})$.

\item Let
\[
\bar{y}_n \equiv \frac{1}{n}(y_1 + y_2 + \cdots + y_n).
\]
Show that
\[
\Var(\sqrt{n}\,\bar{y}_n) = \gamma_0
+ 2\!\left[\Bigl(1 - \frac{1}{n}\Bigr)\gamma_1
+ \Bigl(1 - \frac{2}{n}\Bigr)\gamma_2\right].
\]

\item Because $\{y_t\}$ is not i.i.d., the Lindeberg--Levy CLT is not applicable.
However, it will be shown in Chapter~6 that $\sqrt{n}\,\bar{y}_n$ converges in distribution
to a normal random variable. What is the mean of the limiting distribution
(i.e., $\avar(\bar{y}_n)$)?
\textbf{Hint:} In Lemma~2.1, set $z_n = \sqrt{n}\,\bar{y}_n$, and
$\lim_{n\to\infty}(1 - q/n) = 1$ for any fixed $q$.
\end{enumerate}

\exercise[11] (Optional, Breusch--Godfrey test for serial correlation)\quad
In this question, we
prove that the modified Box--Pierce $Q$ is asymptotically equivalent to $pF$ from
auxiliary regression \eqref{eq:2.10.21}, where $p$ is the number of lags and $F$ is the $F$-ratio for the hypothesis that the coefficients of $e_{t-1}, e_{t-2}, \ldots, e_{t-p}$ are all zero.

First we establish the notation. Let
\[
\underset{(n\times K)}{\mathbf{X}}
= \begin{bmatrix} \mathbf{x}_1' \\ \vdots \\ \mathbf{x}_n' \end{bmatrix}, \quad
\underset{(n\times p)}{\mathbf{E}}
= \begin{bmatrix}
0 & 0 & \cdots & 0 \\
e_1 & 0 & \cdots & \vdots \\
e_2 & e_1 & & 0 \\
e_3 & e_2 & & e_1 \\
\vdots & \vdots & & \vdots \\
e_{n-1} & e_{n-2} & \cdots & e_{n-p}
\end{bmatrix}, \quad
\underset{(n\times 1)}{\mathbf{e}}
= \begin{bmatrix} e_1 \\ \vdots \\ e_n \end{bmatrix},
\]
\[
\underset{((K+p)\times(K+p))}{\widehat{\mathbf{B}}}
= \begin{bmatrix}
\frac{1}{n}\mathbf{X}'\mathbf{X} & \frac{1}{n}\mathbf{X}'\mathbf{E} \\[4pt]
\frac{1}{n}\mathbf{E}'\mathbf{X} & \frac{1}{n}\mathbf{E}'\mathbf{E}
\end{bmatrix}, \quad
\widehat{\mathbf{B}}^{-1} = \begin{bmatrix}
\widehat{\mathbf{B}}^{11} & \widehat{\mathbf{B}}^{12} \\
\widehat{\mathbf{B}}^{21} & \widehat{\mathbf{B}}^{22}
\end{bmatrix}, \quad
\underset{(p\times 1)}{\hat{\bm{\gamma}}}
= \begin{bmatrix} \hat{\gamma}_1 \\ \vdots \\ \hat{\gamma}_p \end{bmatrix},
\]
where $e_t$ is the residual from the original regression
$y_t = \mathbf{x}_t'\bm{\beta} + \varepsilon_t$ and $\hat{\gamma}_j$ is the sample $j$-th order
autocovariance, defined in \eqref{eq:2.10.8}, calculated from the residuals.
\begin{enumerate}[label=(\alph*)]
\item The auxiliary regression \eqref{eq:2.10.21} has $K + p$ regressors. Let
$\hat{\bm{\alpha}}$ be the vector of the $K + p$ coefficients. Show that
\[
\hat{\bm{\alpha}} = \widehat{\mathbf{B}}^{-1}
\begin{bmatrix} \mathbf{0} \\ {}_{(K\times 1)} \\[2pt] \hat{\bm{\gamma}} \\ {}_{(p\times 1)} \end{bmatrix}.
\]
\textbf{Hint:} Since $\mathbf{e}$ is the vector of residuals from the original regression with
$\mathbf{X}$ as regressors, $\mathbf{X}'\mathbf{e} = \mathbf{0}$.

\item Show:
\[
\widehat{\mathbf{B}} \pto \mathbf{B}
\equiv \begin{bmatrix}
\bm{\Sigma}_{\mathbf{xx}} & \mathbf{H} \\
\mathbf{H}' & \sigma^2\mathbf{I}_p
\end{bmatrix},
\]
where
\[
\underset{(K\times p)}{\mathbf{H}}
\equiv \begin{bmatrix}
\E(\mathbf{x}_t\cdot\varepsilon_{t-1}) & \cdots &
\E(\mathbf{x}_t\cdot\varepsilon_{t-p})
\end{bmatrix}.
\]
\textbf{Hint:} The $j$-th column of $\frac{1}{n}\mathbf{X}'\mathbf{E}$ is
$\frac{1}{n}\sum_{t=j+1}^{n}\mathbf{x}_t\cdot e_{t-j}$. Use (2.3.9) to show
that it converges in probability to $\E(\mathbf{x}_t\cdot\varepsilon_{t-j})$.

\item (very easy)\quad Show: $\hat{\bm{\alpha}} \pto \mathbf{0}$.

\item Let $SSR$ here be the sum of squared residuals from the auxiliary, \emph{not} the
original, regression. Show:
\[
\frac{SSR}{n - K - p} \pto \sigma^2,
\]
where $\sigma^2$ is the variance of $\varepsilon_t$, the error term from the original regression.
\textbf{Hint:}
$[\mathbf{X} \;\vdots\; \mathbf{E}]\,\hat{\bm{\alpha}} = \mathbf{E}\hat{\bm{\gamma}}$.
So $SSR = (\mathbf{e} - \mathbf{E}\hat{\bm{\gamma}})'(\mathbf{e} - \mathbf{E}\hat{\bm{\gamma}})$.
$\frac{1}{n}\mathbf{E}'\mathbf{e} = \hat{\bm{\gamma}}$,
$\frac{1}{n}\mathbf{E}'\mathbf{E} \pto \sigma^2\mathbf{I}_p$.

\item Show:
\[
pF = \frac{n\hat{\bm{\gamma}}'\widehat{\mathbf{B}}^{22}\hat{\bm{\gamma}}}
{SSR/(n - K - p)}.
\]
\textbf{Hint:} Apply the formula for the $F$-ratio in (1.4.9) to the auxiliary regression.

\item Use the formula for the partitioned inverse (see (A.10) of Appendix~A) to
show
\[
\widehat{\mathbf{B}}^{22} = \Bigl[\frac{1}{n}\mathbf{E}'\mathbf{E}
- \Bigl(\frac{1}{n}\mathbf{E}'\mathbf{X}\Bigr)\mathbf{S}_{\mathbf{xx}}^{-1}
\Bigl(\frac{1}{n}\mathbf{X}'\mathbf{E}\Bigr)\Bigr]^{-1}.
\]

\item Let $\hat{\bm{\rho}}$ and $\widehat{\mathbf{\Phi}}$ be as in \eqref{eq:2.10.8}
and \eqref{eq:2.10.19}. Show that the modified Box--Pierce $Q$ defined in \eqref{eq:2.10.20}
is asymptotically equivalent to $pF$ (i.e., the
difference between the two converges to zero in probability).
\textbf{Hint:} Show that both $pF$ and the modified $Q$ are asymptotically equivalent to
\[
n\hat{\bm{\gamma}}'(\mathbf{I}_p - \mathbf{\Phi})^{-1}\hat{\bm{\gamma}}/\sigma^2,
\]
where $\mathbf{\Phi}$ is defined in \eqref{eq:2.10.18}.
\end{enumerate}

\exercise[12] (Optional, Chow test for structural change in large samples)\quad
Consider applying the Chow test for structural change to the regression model with conditional
homoskedasticity described in Proposition~2.5. Since the issue of structural
change arises mostly in time-series models, we use ``$t$'' for the subscript. We
generalize Assumption~2.1 by allowing the coefficient vector to change at the
break date $r$:
\[
y_t = \begin{cases}
\mathbf{x}_t'\bm{\beta}_1 + \varepsilon_t & \text{for } t = 1, 2, \ldots, r, \\
\mathbf{x}_t'\bm{\beta}_2 + \varepsilon_t & \text{for } t = r+1, r+2, \ldots, n.
\end{cases}
\]
We assume that the break date $r$ is known. (When the break date is unknown,
the situation is more complicated and has been an object of recent research.
See Stock (1994, Section~5) for a survey.) The null hypothesis to test is that
$\bm{\beta}_1 = \bm{\beta}_2$. This is a set of $K$ restrictions. Let $SSR_1$ be the sum of squared
residuals from the first period ($t = 1, 2, \ldots, r$), $SSR_2$ be the SSR from the
second period ($t = r+1, r+2, \ldots, n$), and $SSR_R$ be the SSR from the entire
sample under the constraint $\bm{\beta}_1 = \bm{\beta}_2$. From Chapter~1, the Chow statistic is
defined as
\[
F = \frac{[SSR_R - (SSR_1 + SSR_2)]/K}
{(SSR_1 + SSR_2)/(n - 2K)}.
\]
Let $\mathbf{b}_1$ be the OLS estimate of $\bm{\beta}_1$ obtained from the first period and
$\mathbf{b}_2$ be the OLS estimate from the second period.
\begin{enumerate}[label=(\alph*)]
\item Show that $KF$ is numerically equal to
\[
\sqrt{n}(\mathbf{b}_1 - \mathbf{b}_2)'\!
\left[\Bigl(\frac{\frac{1}{n}\sum_{t=1}^{r}\mathbf{x}_t\mathbf{x}_t'}{s^2}\Bigr)^{-1}
+ \Bigl(\frac{\frac{1}{n}\sum_{t=r+1}^{n}\mathbf{x}_t\mathbf{x}_t'}{s^2}\Bigr)^{-1}
\right]^{-1}\!\sqrt{n}(\mathbf{b}_1 - \mathbf{b}_2),
\]
where $s^2 = (SSR_1 + SSR_2)/(n - 2K)$. \textbf{Hint:} The $n$ equations can be written
in matrix form as $\mathbf{y} = \mathbf{X}\bm{\beta} + \bm{\varepsilon}$, where
\[
\underset{(n\times 2K)}{\mathbf{X}}
= \begin{bmatrix} \mathbf{X}_1 & \mathbf{0} \\ \mathbf{0} & \mathbf{X}_2 \end{bmatrix}, \quad
\underset{(r\times K)}{\mathbf{X}_1}
= \begin{bmatrix} \mathbf{x}_1' \\ \vdots \\ \mathbf{x}_r' \end{bmatrix}, \quad
\underset{((n-r)\times K)}{\mathbf{X}_2}
= \begin{bmatrix} \mathbf{x}_{r+1}' \\ \vdots \\ \mathbf{x}_n' \end{bmatrix}, \quad
\underset{(2K\times 1)}{\bm{\beta}}
= \begin{bmatrix} \bm{\beta}_1 \\ \bm{\beta}_2 \end{bmatrix}.
\]
The null that $\bm{\beta}_1 = \bm{\beta}_2$ can be written as
$\mathbf{R}\bm{\beta} = \mathbf{r}$ where
$\mathbf{R} = [\mathbf{I}_K \;\vdots\; {-\mathbf{I}_K}]$ and $\mathbf{r} = \mathbf{0}$.
Use Proposition~1.4.

\item Let $\lambda \equiv r/n$. Show that, as $n \to \infty$ with $\lambda$ fixed,
\[
\frac{1}{n}\sum_{t=1}^{r}\mathbf{x}_t\mathbf{x}_t' \pto
\lambda\,\bm{\Sigma}_{\mathbf{xx}}, \quad
\frac{1}{n}\sum_{t=r+1}^{n}\mathbf{x}_t\mathbf{x}_t' \pto
(1-\lambda)\,\bm{\Sigma}_{\mathbf{xx}},
\]
where $\bm{\Sigma}_{\mathbf{xx}} = \E(\mathbf{x}_t\mathbf{x}_t')$.

\item Show that, as $n \to \infty$ with $\lambda$ fixed,
\[
\frac{1}{\sqrt{n}}\sum_{t=1}^{r}\mathbf{x}_t\cdot\varepsilon_t
\dto N(\mathbf{0},\;\lambda\sigma^2\bm{\Sigma}_{\mathbf{xx}}), \quad
\frac{1}{\sqrt{n}}\sum_{t=r+1}^{n}\mathbf{x}_t\cdot\varepsilon_t
\dto N(\mathbf{0},\;(1-\lambda)\sigma^2\bm{\Sigma}_{\mathbf{xx}}).
\]
It can be shown (see, e.g., Stock, 1994, Section~5) that
$\frac{1}{\sqrt{n}}\sum_{t=1}^{r}\mathbf{x}_t\varepsilon_t$ and
$\frac{1}{\sqrt{n}}\sum_{t=r+1}^{n}\mathbf{x}_t\varepsilon_t$
are asymptotically uncorrelated. So
\[
\begin{bmatrix}
\frac{1}{\sqrt{n}}\sum_{t=1}^{r}\mathbf{x}_t\cdot\varepsilon_t \\[6pt]
\frac{1}{\sqrt{n}}\sum_{t=r+1}^{n}\mathbf{x}_t\cdot\varepsilon_t
\end{bmatrix}
\dto N\!\left(\mathbf{0},\;
\begin{bmatrix}
\lambda\sigma^2\bm{\Sigma}_{\mathbf{xx}} & \mathbf{0} \\
\mathbf{0} & (1-\lambda)\sigma^2\bm{\Sigma}_{\mathbf{xx}}
\end{bmatrix}\right).
\]

\item Show that
$\avar(\mathbf{b}_1 - \mathbf{b}_2)
= \lambda^{-1}\sigma^2\bm{\Sigma}_{\mathbf{xx}}^{-1}
+ (1-\lambda)^{-1}\sigma^2\bm{\Sigma}_{\mathbf{xx}}^{-1}$.
\textbf{Hint:} Let
\[
\mathbf{b} \equiv \begin{bmatrix} \mathbf{b}_1 \\ \mathbf{b}_2 \end{bmatrix}.
\]
Show that
\[
\avar(\mathbf{b}) = \begin{bmatrix}
\lambda^{-1}\sigma^2\bm{\Sigma}_{\mathbf{xx}}^{-1} & \mathbf{0} \\
\mathbf{0} & (1-\lambda)^{-1}\sigma^2\bm{\Sigma}_{\mathbf{xx}}^{-1}
\end{bmatrix}.
\]
$\mathbf{b}_1 - \mathbf{b}_2 = [\mathbf{I}_K \;\vdots\; {-\mathbf{I}_K}]\mathbf{b}$.
Use Lemma~2.4(c).

\item Finally, show that $KF \dto \chi^2(K)$ as $n \to \infty$ with $\lambda = r/n$ fixed.
\end{enumerate}

\subsection*{Empirical Exercises}

\exercise[1]
Read the introduction and Sections I--IV of Fama (1975), before doing this
exercise. In the data file \texttt{MISHKIN.ASC}, monthly data are provided on:

\begin{quote}
Column~1: year \\
Column~2: month \\
Column~3: one-month inflation rate (in percent, annual rate; call this \emph{PAI1}) \\
Column~4: three-month inflation rate (in percent, annual rate; call this \emph{PAI3}) \\
Column~5: one-month T-bill rate (in percent, annual rate; call this \emph{TB1}) \\
Column~6: three-month T-bill rate (in percent, annual rate; call this \emph{TB3}) \\
Column~7: CPI for urban consumers, all items (the 1982--1984 average is set to 100; call this \emph{CPI})
\end{quote}

The sample period is February 1950 to December 1990 (491 observations).
The data on \emph{PAI1}, \emph{PAI3}, \emph{TB1}, and \emph{TB3} are the same data used in Mishkin
(1992) and were made available to us by him. The T-bill data were obtained
from the Center for Research in Security Prices (CRSP) at the University of
Chicago. The T-bill rates for the month are as of the last business day of the
previous month (and so can be taken for the interest rates at the beginning of the
month). The construction of \emph{PAI1} and \emph{PAI3} will be described toward the end
of this exercise; for the time being, we will use the inflation derived from \emph{CPI}.

\begin{enumerate}[label=(\alph*)]
\item (Library/internet work)\quad
To check the accuracy of the data in \texttt{MISHKIN.ASC}, find relevant tables from back issues of \emph{Treasury Bulletin} or
\emph{Federal Reserve Bulletin} for hard-copy data on T-bill rates. Or visit the
web sites of the Board of Governors (\texttt{www.bog.frb.fed.us}) or the Treasury Department (\texttt{www.ustreas.gov}) to accomplish the same. Can you
find one-month T-bill rates? [Answer: Probably not.] Are the rates in
\texttt{MISHKIN.ASC} close to those in the relevant tables? Do they appear to be
at the beginning of the month?

\item (Library/internet work)\quad
Find relevant tables from back issues of \emph{Monthly
Labor Review} (or visit the web site of the Bureau of Labor Statistics,
\texttt{www.bls.gov}) to verify that the CPI figures in \texttt{MISHKIN.ASC} are correct.
Verify that the timing of the variable is such that a January CPI observation
is the CPI for the month. Regarding the definition of the CPI, verify the
following. (1)~The CPI is for urban consumers, for all items including food
and housing, and is not seasonally adjusted. (2)~Prices of the component
items of the index are sampled throughout the month. When is the CPI for
the month announced?

\item Is the CPI a fixed-weight index or a variable-weight index? \textbf{Hint:} Dig up
your old intermediate macro textbook; graduate macro textbooks won't do.
\end{enumerate}

The one-month T-bill rate for month $t$ in the data set is for the period from
the beginning of month $t$ to the end of the month (as you just verified). Ideally,
if we had data on the price level at the beginning of each period, we would
calculate the inflation rate for the same period as $(P_{t+1} - P_t)/P_t$, where $P_t$
is the beginning-of-the-period price level. We use CPI for month $t - 1$ for
$P_t$ (i.e., set $P_t = CPI_{t-1}$). Since the CPI component items are collected at
different times during the month, there arises the inevitable misalignment of
the inflation measure and the interest-rate data. Another problem is the timing
of the release of the CPI figure. The efficient market theory assumes that $P_t$ is
known to the market at the beginning of month $t$ when the T-bill rates for the
month are set. However, the CPI for month $t - 1$, which we take to be $P_t$, is
not announced until sometime in the following month (month $t$). Thus we are
assuming that people know the CPI for month $t - 1$ at the beginning of month $t$.

\begin{enumerate}[label=(\alph*),start=4]
\item Reproduce the results in Table~\ref{tab:2.1}. Because the T-bill rate is in percent and
at an annual rate, the inflation rate must be measured in the same unit. Calculate $\pi_{t+1}$, which is to be matched with $\mathit{TB1}_t$ (the one-month T-bill rate
for month $t$), as the continuously compounded rate in percent:
\[
\left[\Bigl(\frac{P_{t+1}}{P_t}\Bigr)^{12} - 1\right] \times 100.
\]

\item Can you reproduce \eqref{eq:2.11.9}, with robust standard errors? What is the
interpretation of the intercept?

\item (Davidson--MacKinnon correction of White)\quad
The finite-sample properties
of the robust $t$-ratio might be improved by the Davidson--MacKinnon adjustment discussed in the text. Calculate robust standard errors using the three
formulas discussed in the text. The first is the degrees of freedom correction of multiplying $\widehat{\mathbf{S}}$ by $n/(n-K)$. The second is to calculate $\widehat{\mathbf{S}}$
by formula (2.5.5) with $d = 1$. The third is (2.5.5) with $d = 2$. These correspond to
TSP's OLSQ options \texttt{HCTYPE = 1, 2, 3}.

\item (Estimation under conditional homoskedasticity)\quad
Test market efficiency
by regressing $\pi_{t+1}$ on a constant and $\mathit{TB1}_t$ \emph{under conditional homoskedasticity} (2.6.1). Compare your results with those in (e) and (f). Which part is
different?

\item (Breusch--Godfrey test)\quad
For the specification in (g), conduct the Breusch--Godfrey test for serial correlation with $p = 12$. (The $nR^2$ statistic should
be about 27.0.) Let $e_t$ ($t = 1, 2, \ldots, n$) be the OLS residual from the
regression in (g). To perform the Breusch--Godfrey test as described in
the text, we need to set $e_t$ ($t = 0, 1, \ldots, -11$) to zero in running the auxiliary
regression for $t = 1, 2, \ldots, n$.

\item (Reconstructing Fama)\quad
What are Fama's own point estimate and standard
error of the nominal interest rate coefficient? Are they identical to your
results? (Fama uses different notation, so you need to translate his results
in our notation.) Why is his estimate of the intercept different from yours?
(One obvious reason is that our data are different, but there is another reason.) Optional: What are the differences between his data and our data?

\item (Seasonal dummies)\quad
The CPI used to calculate the inflation rate is not seasonally adjusted. To take account of seasonal factors while still using seasonally unadjusted data in the regression, define twelve monthly dummies,
$M1, M2, \ldots, M12$, and use them in place of a constant in the regression
in (g). Does this make any difference to your results? What happens if you
include a constant along with the twelve monthly dummies in the regression? Optional: Is there any reason for preferring seasonally unadjusted
data to seasonally adjusted?

\item Estimate the Fama regression for 1/53--7/71 using this better measure of
the one-month inflation rate and compare the results to those you obtained
in (e).

\item Estimate the Fama regression for the post-October 1979 period. Is the nominal interest rate coefficient much lower? (The coefficient should drop to
0.564.)
\end{enumerate}

\exercise[2] (Continuation of the Nerlove exercise of Chapter~1, p.~76)

\begin{enumerate}[label=(\alph*),start=9]
\item For Model~4, carry out White's $nR^2$ test for conditional heteroskedasticity.
($nR^2$ should be 66.546.)

\item (optional)\quad
Estimate Model~4 by OLS. This is \emph{Step~1} of the procedure of
Section~2.8. Carry out \emph{Step~2} by estimating the following auxiliary regression: regress $e_i^2$ on a constant and $1/Q_i$. Verify that \emph{Step~3} is what you did
in part~(h).

\item Calculate White standard errors for Model~4.
\end{enumerate}

\subsection*{Monte Carlo Exercises}

\exercise[1] (Degrees of freedom correction)\quad
In the model of the Monte Carlo exercise of
Chapter~1, we assumed that the error term was normal. Assume instead that $\varepsilon_i$
is uniformly distributed between $-0.5$ and $0.5$.

\begin{enumerate}[label=(\alph*)]
\item Verify that the DGP of the second simulation (where $\mathbf{X}$ differs across simulated samples) satisfies Assumptions~2.1--2.5 and 2.7. (It can be shown that
Assumption~2.2 [ergodic stationarity] is also satisfied. This is an implication of Proposition~6.1(d).)

\item Run the second simulation of the Chapter~1 Monte Carlo exercise with the
uniformly distributed error term. In each replication, calculate the usual
$t$-ratio, as before, and compare it with two critical values. The first is the
same critical value (2.042) implied by the $t(30)$ distribution, and the second
is the critical value (1.96) from $N(0, 1)$. Compute the rejection frequency
for each critical value. Which one is closer to the nominal size of 5\%?
\end{enumerate}

\exercise[2] (Box--Pierce vs.\ Ljung--Box)\quad
We wish to verify the claim that the small sample
properties of the Ljung--Box $Q$ statistic are superior to those of the Box--Pierce
$Q$. Generate a string of 50 i.i.d.\ random numbers with mean 0. (Choose your
favorite distribution.) Taking this string as data, do the following.

\begin{enumerate}
\item Calculate the Box--Pierce $Q$ and the Ljung--Box $Q$ statistics (see \eqref{eq:2.10.4}
and \eqref{eq:2.10.5}) with $p = 4$. (The ensemble mean is 0 by construction. Nevertheless take the sample mean from the series when you compute the autocorrelations.)

\item For each statistic, accept or reject the null of no serial correlation at a significance level of 5\%, assuming that the statistic is distributed as $\chi^2(4)$.
\end{enumerate}

\noindent
Do this a large number of times, each time using a different string of 50 i.i.d.\
random numbers generated afresh. For each $Q$ statistic, calculate the frequency
of rejecting the null. If the finite-sample distribution of the statistic is well
approximated by $\chi^2(4)$, the frequency should be close to 5\%. Which statistic
gives you the frequency closer to the nominal size of 5\%? Do we fail to reject
the null too often if we use the Box--Pierce $Q$?

\end{problemset}

\begin{answers}

\subsection*{Analytical Exercises}

\answer{3}
From (2.3.7) of the text, $\mathbf{b} - \bm{\beta} = \mathbf{S}_{\mathbf{xx}}^{-1}\bar{\mathbf{g}}$.
By Kolmogorov, $\mathbf{S}_{\mathbf{xx}} \pto \E(\mathbf{x}_i\mathbf{x}_i')$ and
$\bar{\mathbf{g}} \pto \mathbf{0}$. So $\mathbf{b}$ is consistent. By the Lindeberg--Levy CLT,
$\sqrt{n}\,\bar{\mathbf{g}} \dto N(\mathbf{0}, \mathbf{S})$.
The rest of the proof should now be routine.

\answer{7}
What remains to be shown is that the LHS of ($**$) equals
$\avar(\mathbf{b}) = \bm{\Sigma}_{\mathbf{xx}}^{-1}\,\mathbf{S}\,\bm{\Sigma}_{\mathbf{xx}}^{-1}$.
\[
\bm{\Sigma}_{\mathbf{xx}} = \plim \frac{1}{n}\mathbf{X}'\mathbf{X}
\]
and
\begin{align*}
\mathbf{S} &= \E(\varepsilon_i^2\,\mathbf{x}_i\mathbf{x}_i') \\
&= \E[\E(\varepsilon_i^2 \mid \mathbf{x}_i)\,\mathbf{x}_i\mathbf{x}_i'] \\
&= \E[(\mathbf{z}_i'\bm{\alpha})\,\mathbf{x}_i\mathbf{x}_i']
\quad \text{(by (2.8.1))} \\
&= \plim \frac{1}{n}\sum_{i=1}^{n}(\mathbf{z}_i'\bm{\alpha})\,\mathbf{x}_i\mathbf{x}_i'
\quad \text{(by ergodic stationarity)} \\
&= \plim \frac{1}{n}\mathbf{X}'\mathbf{V}\mathbf{X}
\quad \text{(by definition (2.8.4) of } \mathbf{V}\text{)}.
\end{align*}

\subsection*{Empirical Exercises}

\answer{1c}
The CPI is probably the most widely used fixed-weight price index.

\answer{1e}
The negative of the intercept $= r$ (the constant ex-ante real interest rate).

\answer{1g}
Point estimates are the same. Only standard errors are different.

\answer{1i}
Looking at Fama's Table~3, his $R_t$ coefficient when the sample period is 1/53 to
7/71 is 0.98 with a s.e.\ of 0.10. Very close, but not exactly the same. There are
two possible explanations for the difference between our estimates and Fama's.
First, in calculating the inflation rate for month $t$,
$\pi_{t+1} = (P_{t+1} - P_t)/P_t$, it is
not clear from Fama that $CPI_{t-1}$ rather than $CPI_t$ was used for $P_t$. Second, the
weight for the CPI at the time of his writing may be for 1958. The weight for
the CPI in our data is for 1982--1984. Our estimate of the intercept ($-0.868$)
differs from Fama's (0.0070) because, first, his dependent variable is the negative of the inflation rate and, second, the inflation rate and the T-bill rate are
monthly rates in Fama.

\answer{1j}
The seasonally adjusted series manufactured by the BLS is a sort of two-sided
moving average. Thus, for example, seasonally adjusted $CPI_t$ depends on seasonally unadjusted values of \emph{future} CPI, which is not in $I_t$. Thus if there is no
reason to take account of seasonal factors in the relationship between the inflation rate and the nominal interest rate, one should use seasonally unadjusted
data. If there is a need to take account of seasonal factors, it is better to include
seasonal dummies in the regression rather than use seasonally adjusted data.
Inclusion of seasonal dummies does not change results in any important way.

\answer{1k}
The $R_t$ coefficient drops to 0.807 from 1.014.

\end{answers}
